% Path pebbling toolkit: shared lemmas used throughout the paper


\subsection{Path pebbling toolkit}
\label{sec:toolkit}


Every bound in this paper --- lower and upper, leaf and non-leaf --- is built
from facts about pebbling on paths. We collect them here so that later
sections can cite them freely.

The potential-function framework introduced first will be used to establish
the lower bound by constructing a quantity that cannot increase under
pebbling moves. The weight-function framework will then provide the upper
bound through valid weight assignments. Finally, the two-armed path bound will
allow us to treat non-leaf targets uniformly in the structural lemmas.


\subsubsection{Potential-function framework}


\begin{lemma}[Path potential bound]
\label{lem:path-lb}

Let $P=(v_0,v_1,\ldots,v_k)$ be a path and let $t=v_k$.
Define

\[
\Psi(p)=\sum_{i=0}^{k}p(v_i)2^{-(k-i)}.
\]

Then $\Psi(p)$ is non-increasing under every pebbling move on $P$.

In particular, if

\[
p(v_0)=2^k-1
\]

and

\[
p(v_i)=0\quad(i>0),
\]

then $t$ cannot be pebbled.

\end{lemma}


\begin{proof}

A move $v_i\to v_{i+1}$ satisfies

\[
\Delta\Psi
=
-2\cdot2^{-(k-i)}
+
2^{-(k-i-1)}
=0 .
\]

A move $v_{i+1}\to v_i$ satisfies

\[
\Delta\Psi
=
-2\cdot2^{-(k-i-1)}
+
2^{-(k-i)}
<0 .
\]

Hence $\Psi$ never increases.

For the critical configuration,

\[
\Psi(p)
=
(2^k-1)2^{-k}
<1 .
\]

A pebble on $t$ contributes exactly

\[
2^0=1
\]

to the potential, so $t$ cannot be reached.

\end{proof}



\subsubsection{Weight-function framework}


We use the weight-function framework of
Hurlbert~\cite{Hurlbert2017},
Cranston et al.~\cite{Cranston2017},
and Yang--Yerger--Zhou~\cite{YYZ2024}.


Fix a target $r$.

A weight function is

\[
w:V(G)\rightarrow\mathbb R_{\ge0}
\]

with

\[
w(r)=0
\]

and

\[
w(v)>0\quad(v\neq r).
\]


It is valid if every unsolvable configuration satisfies

\[
w(p)\le w(1_G).
\]


\begin{lemma}[Weight-function bound
{\cite[Lemma~1]{YYZ2024}}]
\label{lem:wf-bound}

If $w$ is valid and

\[
m=\min_{v\neq r}w(v),
\]

then

\[
\pi(G,r)
\le
\left\lfloor\frac{w(1_G)}m\right\rfloor+1 .
\]

\end{lemma}



\begin{lemma}[Path weights
{\cite[Lemma~2]{YYZ2024}}]
\label{lem:yyz-path}

Let

\[
P=(v_0,\ldots,v_{k-1},r)
\]

be a path with

\[
w(v_i)=2^i .
\]

Then

\[
p\text{ is }r\text{-solvable on }P
\Longleftrightarrow
w(p)\ge2^k .
\]

\end{lemma}



\begin{lemma}[Two-armed path bound]
\label{lem:two-arm}

Let

\[
Y=(y_0,y_1,\ldots,y_D)
\]

be a path and fix

\[
j\in\{0,\ldots,D\}.
\]

Then

\[
\pi(Y,y_j)\le2^j+2^{D-j}-1 .
\]

\end{lemma}


\begin{proof}

Split $Y$ into the two paths

\[
Y_L=(y_0,\ldots,y_j),
\]

and

\[
Y_R=(y_j,\ldots,y_D).
\]

The two path weight functions give

\[
w_L(p)\le2^j-1
\]

and

\[
w_R(p)\le2^{D-j}-1
\]

for every unsolvable configuration.

Therefore,

\[
|p|
\le
w_L(p)+w_R(p)
\]

and hence

\[
|p|
\le
2^j+2^{D-j}-2 .
\]

Thus every configuration with

\[
2^j+2^{D-j}-1
\]

pebbles reaches $y_j$.

\end{proof}